% TOP_PROBLEM: {"id": 30007019, "title": "Canonical minimal models in numerical dimension one", "category_id": 5, "category": {"id": 5, "name": "algebraic_geometry", "display_name": "Algebraic Geometry", "description": "Geometric objects defined by polynomial equations.", "slug": "algebraic-geometry", "order_index": 5, "created_at": "2026-07-31T15:26:25.671Z"}, "status": "open"}
% FIRST_POSTED: 2026-10-06
% ENTRY_KIND: statement_only
% !TeX program = lualatex
\documentclass[11pt]{article}
\usepackage[margin=25mm]{geometry}
\usepackage{fontspec}
\setmainfont{Latin Modern Roman}
\usepackage{amsmath,amssymb,mathtools}
\usepackage{unicode-math}
\setmathfont{Latin Modern Math}
\usepackage{xurl,hyperref}
\hypersetup{colorlinks=true,urlcolor=blue}
\setcounter{secnumdepth}{0}
\setlength{\parindent}{0pt}
\setlength{\parskip}{5pt}
\setlength{\emergencystretch}{3em}
\newcommand{\R}{\mathbb R}
\newcommand{\N}{\mathbb N}
\newcommand{\Z}{\mathbb Z}
\newcommand{\Q}{\mathbb Q}
\newcommand{\C}{\mathbb C}
\newcommand{\D}{\mathbb D}
\newcommand{\F}{\mathbb F}
\newcommand{\sref}[2]{\hyperlink{source:#1:#2}{[S#2]}}
\newcommand{\cataloguescope}[1]{#1}
\begin{document}
\section*{Canonical minimal models in numerical dimension one}\label{30007019}
\subsection{Definitions and mathematical statement}
\paragraph{Statement recorded in the supplied source.}
If a smooth connected projective complex n-fold, $n\geq 3$, has pseudo-effective $K_X$ with Nakayama numerical Kodaira dimension $\kappa _\sigma (K_X)=1$, must a finite sequence of K-negative divisorial contractions and flips reach a projective Q-factorial terminal model Y with nef $K_Y$? On a common resolution must $p*K_X=q*K_Y+E$ hold with E effective q-exceptional and containing every contracted X-prime in its support?
\subsection{Short English statement}
If a smooth connected projective complex n-fold, $n\geq 3$, has pseudo-effective $K_X$ with Nakayama numerical Kodaira dimension $\kappa _\sigma (K_X)=1$, must a finite sequence of K-negative divisorial contractions and flips reach a projective Q-factorial terminal model Y with nef $K_Y$? On a common resolution must $p*K_X=q*K_Y+E$ hold with E effective q-exceptional and containing every contracted X-prime in its support?
\subsection{Sources}
\begin{itemize}
\item[S1] Original question source or mathematical context.
\url{https://arxiv.org/abs/2406.18233}.
\item[S2] Original question source or mathematical context.
\url{https://arxiv.org/abs/2601.15138}.
\item[S3] Original question source or mathematical context.
\url{https://arxiv.org/abs/2604.03879}.
\item[S4] Original question source or mathematical context.
\url{https://arxiv.org/html/2501.04232v1}.
\item[S5] Original question source or mathematical context.
\url{https://doi.org/10.4171/owr/2014/10}.
\item[S6] Original question source or mathematical context.
\url{https://publications.mfo.de/bitstream/handle/mfo/3399/OWR_2014_10.pdf?isAllowed=y&sequence=1}.
\item[S7] Original question source or mathematical context.
\url{https://www.cambridge.org/core/journals/nagoya-mathematical-journal/article/two-nonvanishing-results-concerning-the-anticanonical-bundle/1E68F2B88B818A7CBD47B96B7D700548}.
\item[S8] OpenAI, manuscript citation page. Source of the formulation motivating this local problem.
\url{https://github.com/openai/math/blob/adc7f1241b42e322a6451854ab7e4b4c146bf78a/preprints/Minimal-models-in-numerical-dimension-one-September-24-2026/README.md}.
\end{itemize}
\end{document}
